% =====================================================================
\section{Conclusion}\label{sec:conclusion}
% =====================================================================

We set out to see what happens when Rust's ownership mechanisms meet a dependently typed kernel and a hygienic macro system.

\emph{RQ1: kinds and elimination.}
Function kinds and dependent elimination are not independent: a recursor runs its case functions as often as the data requires, so a case that may run more than once must not consume what it captures, and a recursive field is consumed when its induction hypothesis is computed.
The kernel enforces these rules; for $\lambda^{\mathrm{aff}}$, a non-dependent core calculus with function kinds and an iterator that models a repeatable case, we mechanized that well-owned programs never consume a resource twice, while the dependent calculus is not mechanized.

\emph{RQ2: who enforces what.}
The kernel, which is the logical trusted base, enforces affinity; the MIR borrow checker, outside the kernel, enforces borrowing; every definition passes both, and macros cannot avoid either because they run before elaboration and see neither types nor ownership.
The two checkers check different things and are not required to agree; on the repository's programs the kernel never rejected a definition that MIR accepted, and on the violation corpus MIR independently rejected all but one class of the kernel's ownership violations.

\emph{RQ3: does it work?}
Complete programs that use all three features together can be written, checked and compiled.
On our twelve properties, Idris~2 statically rejects every violation that LRL rejects but one (it accepts a proof forged with \lstinline|believe_me|; the program then fails at run time in the division the proof was meant to protect), and two that LRL accepts; LRL's design differs in expressing resources through Rust-style closure kinds and borrows with lifetimes.
LRL does not yet turn ownership into efficiency, which Lean~4 (through reference counting), Idris~2 (with linear arrays) and Rust do.

The most immediate directions for future work follow from the limitations: operations that read and write through references, a reuse analysis that exploits affinity for in-place update, and a mechanization that covers recursors and dependent types.

\section*{Data and code availability}

The implementation, the case studies, the comparison programs, the measurement scripts and the Lean development are available at \url{https://github.com/Kendiukhov/LeanRustLisp} under the MIT licence.
All results in this paper refer to commit \texttt{\PinnedCommit}.
The mechanization builds with \texttt{lake build LRL} in \texttt{mechanization/}; the compiler and its \NumTestsPassed{} tests build and run with \texttt{cargo test --all}; each generated report names the script that regenerates it (for the measurements, \texttt{bench/\allowbreak results/\allowbreak summary.md}, which also points to the raw data files).

\section*{Declaration of competing interest}

The author declares no competing financial interests or personal relationships that could have influenced the work reported in this paper.
